\chapter[Sermon 92. It Is Furthermore Declared, That the Hope of the Everlasting and Blessed Felicity and Glory to Be Certain and Undoubted.]{It Is Furthermore Declared, That the Hope of the Everlasting and Blessed Felicity and Glory to Be Certain and Undoubted.}
\label{sermon:092}
\markright{Sermon 92. It Is Furthermore Declared, That the Hope of the ...}

I am Alpha and Omega, the beginning and the end. I will give to the thirsty person from the well of the water of life freely. He who overcomes shall inherit all things. I will be his God, and he shall be my sun. But the fearful and unbelieving, and the abominable, and murderers, and adulterers, and sorcerers, and idolaters and liars shall have their part in the lake that burns with fire and brimstone, which is the second death.

To all the former comes now the sixth testimony of the certainty of the true felicity of the faithful, taken from the very nature of God. For he proclaims of himself and says, “I am Alpha and Omega.” And immediately by explanation, “the beginning and the end.” He took this from Isaiah, with whom the Lord says repeatedly, “I am first and last.” And here let no man imagine that God is first in order, referring the beginning to consequences, as though he had a beginning; or that he is called the last or end, as though he should ever have an end; but the contrary rather is to be understood, namely that God has no beginning, no end, but is everlasting, from whom all things have their being, and by whose decree all things have an end; where he himself endures forever, and his years never fade, just as the prophet says in another place, and the Apostle also. And since he is eternal, without beginning and without end, who always lives, and quickens and preserves all things that live, how, I pray you, should he not quicken the faithful? So certain therefore is the life, salvation, and felicity of the faithful, as it is certain that God is life, and indeed everlasting life. For he is everlasting, and the life of the faithful. I have spoken about the phrase of speech, “I am Alpha and Omega,” in the first Chapter and third Sermon.

The seventh witness of our assured salvation comes from the truth of God and his promises, and it shares a certain connection with the previous ones. For that which God has promised, he can certainly fulfill without difficulty. He has promised a blessed life, and therefore he will assuredly grant it to the faithful. He cites God’s promise at this time, bringing in God speaking to John and to us with these words: “To him who thirsts I will give of the well of living water.” That is to say, I, who am life and eternal life itself, will give the faithful to drink the water of life, that is to say, I will quicken him, preserve him in life, deliver him from death and all evils, and reward him with all heavenly gifts. Who can doubt the truthfulness of the one who promises, especially since this place or this promise is found in more than one location? David sings plainly in Psalm 36: “Your mercy, O Lord, reaches to the heavens; and your faithfulness to the clouds; your righteousness is like the strong mountains, your judgments like the great deep. You, O Lord, save both man and beast. How excellent is your mercy, O God! And the children of men shall put their trust under the shadow of your wings. They shall be satisfied with the fullness of your house; and you shall give them drink out of the river of your pleasures. For with you is the well of life, and in your light shall we see light.” Many of these things are found in the Prophets, and are explained by our Savior himself in chapters 4 and 7 of the Gospel of John, where he shows that he gives water and wholesome drink to the faithful, which at length will spring up into eternal life. It is therefore certain that the faithful are quickened by Christ; and therefore, the blessed life of the faithful is, and will be, most assured and certain, as promised by so many explicit promises of God. We discussed this water of life somewhat in chapter 7 of this book toward the end, and we will have certain clear matters in the beginning of chapter 22.

But in the meantime and by the way, he shows and declares unto us, after the apostolic manner, who willingly and often declare unto us the manner of our salvation, how eternal life is communicated to us, to wit, freely, graciously, which notwithstanding, for the doubtfulness of speech or understanding of words, we do not properly express the force of the Greek word. They are justified, says the Apostle in Romans 3:24, freely through his grace: that is to say, by the mere mercy of God, by no merit of man. For the same Apostle in the same Epistle to Romans, chapter 6, “The reward of sin is death,” he says, “and where on the contrary side he should have set, and the merit of righteousness, eternal life, for this member he places rather, and the gift of God is life everlasting.” And he adds immediately, through Christ Jesus our Lord. Therefore John rightly says that eternal life happens to the faithful freely: that is, by the very grace of God, through the merit of Christ, and by no desert of man. For if we could by our works and righteousness deserve eternal life, then Christ had died in vain, to no purpose. There was no cause why he should die, saying we might of ourselves have been saved. There is no effect, nor merit of Christ’s passion, such an effect as it is in very deed, that by the blood of Christ alone we are purified. For if there were or had been another means of salvation, Christ needed not to have been incarnated and have suffered. And that this term ought after this way and manner to be explained, many other places of Scripture prove. In Matthew 10:8, the Lord says, “Freely you have received, freely give.” The Lord will not have his Apostles to receive any recompense for the gift of healing. But speaking of the ministry, he says, “The workman is worthy his hire.” In John 15, the Lord says, “They have hated me without cause, doubtless without my desert, or undeserved of my part.” In 2 Corinthians 11, the Apostle says that he preached the gospel freely, without taking reward or recompense therefore. And in 2 Thessalonians 3, he says, “Neither have I taken bread of any man for naught, in short, where John says that life is given to the faithful, he attributes all things of our salvation to the grace of God, and the merit of Christ’s passion, and plucks it from man’s merits. And the same affirms Isaiah also in chapter 55, rebuking foolish men, spending their money about things of naught. Here ought therefore to cease the fairs of indulgences and pardons, and holy things in the church. Let the Pelagians keep silence.

However, lest anyone, by the free preaching of the grace and merit of Christ, against the deservingness of humans, should gather that the blessed life falls to idle people, sleepers, and those ceasing from all good works: and that God alone works, and we work nothing: but only suffer the operation of God in us, and for that cause nothing is required of us, he anticipates this. First, the Lord says that he will give to those who thirst to drink of the water of life. Therefore, faith and a fervent desire for godly things are required of us; not that faith is prior [?], but that it is given by God. For by thirst, to signify the faithful desire of a godly man, the Lord himself is the author in Matthew 5, pronouncing blessed those who hunger and thirst for righteousness. And also in John 6, the Lord himself understands by drinking to believe. Therefore, faith is required of us, that is, that we should thirst for the water of life. This very thing also the Lord grants by his Spirit and word, as we have declared elsewhere. And he says how he who is freely justified must also fight; not only fight, but must overcome. Therefore, the duties of charity are required, of which there is frequent mention in chapters 2 and 3 of this book, wherein is made most frequent mention of this fight and victory. And God will acknowledge those who labor thus valiantly for his children; to them he will show himself a father, and take them as the heirs of all their fathers’ possessions. They are bastardly children, who, being idle, boast of faith, praise God with their mouth and words, and deceive him with their deeds. You see therefore that both must be preached in the church: that we are justified and beautified freely; and so, being justified, must work good works, to which, notwithstanding, as to their merits, they ascribe not salvation, but to the mere grace of God through Christ.

Consequently, and on the contrary, he lists those excluded from the fellowship of blessed life and the blessed, compiling a register of sins and wicked men, as he has done at the end of chapter 9, verses 21 and 22. And as the Apostle has, in a manner, recited to the Corinthians.

\index{Repentance}And we suppose that in John’s time these sins were most common, nor sufficiently known as was fitting. Many also at this day judge them more lightly than true godliness permits. And we doubt not but that in this record, which is comprised in eight kinds or members, are contained all other like sins and wickednesses. But we understand that hell fire is assuredly due to them for their sins committed, who have no faith at all, neither can by any means be persuaded to repent and turn unto God. For in the first Epistle to the Corinthians, chapter 6, he says, “You were such, but you have been cleansed by the blood of Christ and by the Spirit of our God.” Therefore, if we have been such at any time, let us repent; or, if we have fallen into these sins again, let us rise up and turn to the Lord, who calls to him sinners and promises pardon and grace. But woe to the incurable, walking always and without repentance in the way of iniquity.

\index{Jerome, Saint}And we shall consider separately eight aspects of this register. First are listed the fearful. But the Lord himself was afraid, and even quaked for fear of death: the saints of God have feared also, and often fled for fear: yet are they not for this cause condemned in the Scriptures. Therefore another fear is meant, namely that immoderate fear, by which, compelled, we do for fear of men, that thing which God has prohibited: and we ourselves convict in our own consciences, understand that we sin in so doing: or when, through carnal fear, we leave undone that thing which God has commanded us: briefly, when we more fear men, as princes or leige fellows, or enemies, or any other men whatever they be, than our Lord God himself. And therefore the Lord himself in the gospel said: fear not them which kill the body, and cannot kill the soul, \&c. Matthew 10. The same in another place says: he that denies me before that aduouterous generation, I will deny him also before my father in heaven. Doubtless it is a foul shame to fear more a most wicked man than most holy God. But men offend in this regard at these days most grievously. For some attribute so much to wicked and cruel persecutors that even for them they will command to pervert the preaching of the Gospel, or to keep silence altogether. There is one who will set more by the King, Prince, Earl, Baron, Citizen, or plowman, Bishop or Abbot, or some flattering friar, or vile massmongering priest, and will fain and dissemble for his favor, rather than he will freely confess the truth, and fear and glorify God, who is to be feared only. Unto them says Isaiah: say not conspiracy, and be not afraid of terror of the enemies, neither be you discouraged. But rather sanctify the Lord of hosts: let him be your terror, let him be your fear. He shall be the sanctuary, and stumbling stone, and the rest in the eighth chapter of Isaiah. For unless we put away this vain and wicked fear, and go about to furnish up the Lord’s work valiantly, constantly, and without fear, we shall surely be cast down to hell. Let fearful men think hereof, and call upon the Lord, and take unto them the spirit of strength, and of wise and godly boldness: and do the work of the Lord not negligently, but diligently, valiantly, and constantly. He is greater that is in us, says Jerome in his canonical writings, than is he that is in the world.

Unbelievers are not weak in the faith, modest, and fearing God; but such as do not believe God’s word, its promises, commands, and threats, neither follow God nor his Christ. Rather, they follow strange gods, prefer to believe fables, and have withdrawn their hearts from God. And of these there is a great multitude at this day, who nevertheless have all in their mouth that they believe God and his word, but they do not believe the preachers, thinking truly that their incredulity is thus sufficiently excused. But where the preachers show nothing else but the word of God, they cannot but restrain God’s word while they despise the sermons of the preachers.

\index{Hebrews, Epistle to the}In the third place follows that the torments of Hell are due to the abominable and detestable. For “abomination” signifies abomination and stench. He notes therefore abominable and detestable men, to whom all religion is a mockery, who deride God and his word, and blaspheme all holy things, the children of Belial, uncurable, and scornful. These, although they know the truth, yet they know it to their own condemnation, seeing they contemn it, and with dogs and hogs return to their vomit and wallowing in the mire. Whom also the Apostles have noted: Peter in the second epistle, chapter 2, verses 1-3; Paul in 3 Timothy; and 12 Hebrews; Jude Thaddeus throughout the chiefest part of his epistle; John himself about the end of chapter 22, reciting in a manner the same register, calls them dogs. And would God we lacked examples at this day of abominable men and such kind of dogs. But there is no cause why we should marvel hereat, considering that we live in the time of all others, the most corrupt of Noah and Lot, as Matthew 24.

\index{Rome}Of homicides there are sundry kinds. For we kill with the heart, mouth, and works. You may see the expositors of the ten commandments on this, chiefly Dr. Musculus. But I think the world has never known a more notable, more cruel, and more shameless murderer, indeed a parricide most certainly, according to the word of Christ in John 8:44, the firstborn son of the Devil, than the Bishop of Rome. For he, in a manner, at all times, for these five hundred years and more, has blown the trumpet to all the grievous wars of Europe or Christendom; and again has granted to murderers, especially those warring for the See of Rome, most large and ample pardons, and promised heaven to them that die in that warfare. All the which, were it not for the great mercy of God, would have destroyed both body and soul.

\index{John, Saint}\index{Papacy}Then Saint John recounts fornicators. And he names the lowest kind, so that we should understand the higher and more vile ones, such as rape, adultery, incest, and Sodomitical acts; nor should we exclude gluttony, drunkenness, and all kinds of riotousness and the fostering of voluptuousness. Wherein we undoubtedly see that Saint Paul, under the term fornicators, comprehends all filthiness and riot. But in our days, whoring is so common that every most shameful fornicator is admitted to the altar; a married priest who keeps holy matrimony is expelled from it. For this we may thank Syricius and other Popes, whom the Apostle has severely noted in the first letter to Timothy, 4.

Note 92.12: Sorcerers are discussed in chapter 9 of this book. John [? describes] them, and by this he means magic, enchanters, diviners, witches, and those who use devilish arts to influence affections. The Latin authors also understand them to be those who administer poison.

Idolaters are those who worship idols. And it is astonishing that the papists of today deny that they are idolaters. For what else is an idol but a shape or image made of some visible material, representing the form of God or a saint, but without spirit? An idol, therefore, is an image of wood, stone, or metal, representing the shape of God the Father, of God the Son, or of Saint Peter, and so on. David describes an idol, and says: “The idols of the heathen are silver and gold, the work of human hands. They have mouths but do not speak; they have eyes but do not see.” Psalm 113. And I would fain know what difference the idols of the papists have from these? Concerning worshipping them, they cannot deny that they worship those idols of wood and clay. For they attribute to them holy names, and even the sacred name of God, which is communicated to no other, saying: “This [pointing to stone or wood, that is, an idol of wood] is God the Father, this is God the Son, this is Saint Peter.” I tremble in my mind while reporting these things, especially since the Lord himself has said, “Whom will you liken me to?” Isaiah 40. And Saint Paul calls this plainly counterfeiting foolishness, and expressly denies that the godhead is like a stone artificially polished. Acts 17:29. Again, these images, which they call their gods and saints, made with human hands, they bring into the churches, namely a place of worship, and set them upon the altars. To these they go on pilgrimage, fall down before them and worship, kiss them, offer oblations to them, and hang jewels on them. Moreover, they attribute to them also a part of the heavenly doctrine and instruction, saying that the unlearned are taught and admonished by these. And what is worship if this is not? Let them see, therefore, whether they can excuse themselves in this before God and men, and provide rather to save their souls. However, they sweep away all these things as if with one word, and say: “We do not worship the signs, but the things signified.” Then, if the signs were taken away, would they not return to the idols on pilgrimage? Do they not think it done in a manner to God himself that which is done to the idols? Do they not punish an image breaker as a traitor against the divine majesty? For he shall not seem to have cut asunder wood, but to have defiled God himself. Therefore, they acknowledge something more in this wood than wood alone. For you think that some divine thing is hidden therein, and therefore this wood is accounted by you not common wood. Which thing you declare also by sundry tokens otherwise. Moreover, the gentiles excused themselves in the same manner, saying that they worshipped the things, and not the signs. But this seemed not a sufficient excuse unto godly men, as it is to be read in Lactantius and Athanasius in their books against the gentiles. But God has at one word refuted you and said, “Who has required these things at your hands? If any will exhibit to me worship, let him worship according to the prescription of my most holy law. They worship me in vain, teaching the doctrines of men.” These things I have declared somewhat more at large, to the intent that those who will yet hear any reason, and in whom the word and law of God have any place, might know and avoid that gross and mortal sin of idolatry.

And liars include those who are quick of speech, slanderers, gossips, whisperers, deceivers, covetous people, thieves, usurers, bribe-takers, and all kinds of hypocrites and unreliable people. For just as God is truth, so he loves truth, simplicity, steadfastness, and integrity. This vice of lying prevails widely today. For there is the least, or rather no faith at all in the world. The Lord be merciful to us.

Regarding the lake or pond, burning with fire and brimstone, and of the second death, I have spoken before in the nineteenth and twentieth chapters. And elsewhere. He signifies that all these and the like shall be cast down by the Lord into the everlasting fire of Hell. For he puts here a portion as inheritance, as also in Psalm 11: he shall rain fire and brimstone upon the ungodly, and this is part of their cup. And in Matthew 24: he shall put his portion with hypocrites. And we say also, he has not obtained his right; or he is punished as he is worthy. Just as the saints obtain the Kingdom of Heaven by inheritance, so are everlasting torments in place of inheritance to the ungodly. To the Lord, the righteous Judge, be praise and glory. Amen.
